\section{Tooling and build surprises}

\subsection{find\_package rejected the version I pinned}

The first wall was CMake. I wanted to pin the LLVM major version, so I wrote
\texttt{find\_package(LLVM 21 REQUIRED CONFIG)}. It failed, reporting that the
\texttt{LLVMConfig.cmake} advertising version 21.1.8 was considered but not
accepted. LLVM's generated version config does not treat a bare major of 21 as
compatible with the concrete 21.1.8 it ships. Requesting the full patch version
would hard code a number that changes on every apt update. The fix was to take
whatever configuration \texttt{LLVM\_DIR} points at and enforce the pin myself
with a fatal error guard on \texttt{LLVM\_VERSION\_MAJOR}. The pin stays at major
granularity and still fails loudly on the wrong toolchain.

\subsection{The plugin linked, the unit tests did not}

Linking the standalone unit test executable against the LLVM component libraries
surfaced a second issue: \texttt{LLVMSupport} declares a link dependency on
\texttt{zstd::libzstd\_shared}, an imported target that did not exist because I
had not asked for it. The distribution LLVM was built with zstd. The fix was to
resolve the zstd package before LLVM's config is read, with a small shim that
constructs the imported target from the library on disk if the system's zstd
ships no CMake package, so a continuous integration machine stays portable.

\subsection{Spaces in the checkout path broke lit}

The plugin built and ran by hand, but the lit smoke test failed with FileCheck
complaining of too many positional arguments and a library that could not be
loaded. The checkout lives under a Windows Documents folder mounted into WSL,
whose path contains spaces. lit substitutes paths as raw text into a RUN line
that the shell then splits on those spaces. Rather than move the project, I made
quoting a project convention: every RUN line wraps its substitutions in double
quotes. This is strictly more robust, since it also survives a user cloning into
a spaced path, which LLVM's own in tree tests never have to consider.

\subsection{optnone made the whole pipeline a no-op}

The most surprising one. The first time I ran the pipeline on a real kernel, it
reported zero rewrites on functions I knew were full of identities.
\texttt{clang -O0} attaches the \texttt{optnone} attribute to every function, and
the new pass manager honors it by skipping optimization entirely. My passes were
loaded and registered; they were simply never run. The fix was one flag,
\texttt{-Xclang -disable-O0-optnone}, which keeps the naive unoptimized baseline
I wanted while letting \texttt{opt} actually transform it. The harness sets it for
every kernel and records it in the manifest.
